\documentclass{article}
\usepackage{amsmath,amssymb,amsthm,thmtools,enumitem,bm,graphicx,xcolor,microtype}
\usepackage[T1]{fontenc}

\usepackage[style = ext-numeric, sorting = nyt, giveninits = false, maxbibnames = 8, minbibnames = 3, date = year, isbn = false, doi = true, eprint = true, url = true, backref = true]{biblatex}
\addbibresource{../../ZoteroDatabase.bib}
\renewbibmacro{in:}{}
\DeclareFieldFormat{titlecase}{\MakeSentenceCase*{#1}}
\DeclareFieldFormat{titlecase:journaltitle}{{#1}}

\usepackage[unicode]{hyperref}

\theoremstyle{definition}
\newtheorem{dfn}{Definition}[section]
\newtheorem{ntt}[dfn]{Notation}
\newtheorem{ex}[dfn]{Example}
\newtheorem{case}{Case}
\newtheorem{cnd}{Condition}

\theoremstyle{remark}
\newtheorem{rmk}[dfn]{Remark}
\newtheorem*{sroof}{Sketch of Proof}

\theoremstyle{plain}
\newtheorem{thm}[dfn]{Theorem}
\newtheorem*{thm*}{Theorem}
\newtheorem{prop}[dfn]{Proposition}
\newtheorem{lem}[dfn]{Lemma}
\newtheorem{cor}[dfn]{Corollary}
\newtheorem{fct}[dfn]{Fact}
\newtheorem{conj}[dfn]{Conjecture}
\newtheorem{qst}[dfn]{Question}
\newtheorem{clm}{Claim}

\renewcommand{\theHdfn}{\thesection.\arabic{dfn}}
\renewcommand{\theHntt}{\theHdfn}
\renewcommand{\theHex}{\theHdfn}
\renewcommand{\theHrmk}{\theHdfn}
\renewcommand{\theHthm}{\theHdfn}
\renewcommand{\theHprop}{\theHdfn}
\renewcommand{\theHlem}{\theHdfn}
\renewcommand{\theHcor}{\theHdfn}
\renewcommand{\theHfct}{\theHdfn}
\renewcommand{\theHconj}{\theHdfn}
\renewcommand{\theHqst}{\theHdfn}

\renewcommand{\qedsymbol}{$\blacksquare$}

\declaretheoremstyle[headfont = \normalfont\itshape, qed = $\blacksquare_{\mathrm{Claim}}$]{prooc}
\declaretheorem[style = prooc, name = Proof of Claim, numbered = no]{prooc}

\def\sectionautorefname~#1\null{Chapter~#1\null}
\def\subsectionautorefname~#1\null{Section~#1\null}

\title{}
\author{}
\date{}

\begin{document}

\maketitle

\tableofcontents

\section{The infinite spine case}

Throughout this section, $G$ denotes an ordered abelian group.
For $n \geq 2$ and $a \notin nG$, recall that $\mathfrak{s}_n(a)$ is the largest convex subgroup $H \Subset G$ such that $a \notin H + nG$.
If $a \in nG$, we set $\mathfrak{s}_n(a) = \emptyset$, and we regard $\emptyset$ as being smaller than every convex subgroup.
We set $\mathcal{S}_n = \left\{ \mathfrak{s}_n(a) \mid a \in G \right\}$.
The subgroups $\mathfrak{s}_n(a)$ are uniformly definable in $a$ \cite[Lemma 2.1]{Cluckers2011QuaEli}.
Consequently, comparisons and equalities between values of the maps $\mathfrak{s}_n$ are formulas in $L_{\mathrm{oag}}$.

\begin{lem}[{\cite[Lemma 2.8(2) and Remark 2.6]{Okura2026DisOrd}}]\label{lem:BasicPropertiesOfS}
	Let $p$ be a prime, $r \geq 1$, and $a,b \in G$.
	\begin{enumerate}
		\item If $0 \leq j < r$, then
		\[
			\mathfrak{s}_{p^r}(p^j a) = \mathfrak{s}_{p^{r-j}}(a).
		\]
		If $j \geq r$, then $\mathfrak{s}_{p^r}(p^j a) = \emptyset$.
		\item If $\mathfrak{s}_{p^r}(a) \neq \mathfrak{s}_{p^r}(b)$, then
		\[
			\mathfrak{s}_{p^r}(a+b)
			=
			\max \left\{ \mathfrak{s}_{p^r}(a),\mathfrak{s}_{p^r}(b) \right\}.
		\]
	\end{enumerate}
\end{lem}

\begin{lem}\label{lem:RepresentativeBetweenSValues}
	Let $p$ be a prime and let $H<H'$ be nonempty elements of $\mathcal{S}_p$.
	There exists $a\in H'$ such that
	\[
		\mathfrak{s}_p(a)=H
		\quad\text{and}\quad
		\mathfrak{s}_{p^r}(a)<H'
	\]
	for every $r\geq 1$.
\end{lem}

\begin{proof}
	Choose $b\in G$ such that $\mathfrak{s}_p(b)=H$.
	Since $H'>H$, the maximality of $H$ gives $b\in H'+pG$.
	We may therefore write $b=a+pc$ with $a\in H'$ and $c\in G$.
	As $a-b\in pG$, we have $\mathfrak{s}_p(a)=\mathfrak{s}_p(b)=H$.

	Fix $r\geq 1$.
	If $\mathfrak{s}_{p^r}(a)=\emptyset$, there is nothing to prove.
	Otherwise, put $K=\mathfrak{s}_{p^r}(a)$.
	If $K\geq H'$, then $a\in H'\subseteq K$, contrary to $a\notin K+p^rG$.
	Hence, $K<H'$.
\end{proof}

\begin{lem}[Finite coding]\label{lem:FiniteCoding}
	Let $p$ be a prime and
	\[
		K_0<K_1<\dots<K_N
	\]
	be nonempty elements of $\mathcal{S}_p$.
	There exists $x\in G$ such that
	\[
		\mathfrak{s}_{p^{i+1}}(x)=K_i
	\]
	for every $0\leq i\leq N$.
\end{lem}

\begin{proof}
	For each $j<N$, use \autoref{lem:RepresentativeBetweenSValues} to choose $a_j\in G$ such that
	\[
		\mathfrak{s}_p(a_j)=K_j
		\quad\text{and}\quad
		\mathfrak{s}_{p^r}(a_j)<K_{j+1}
	\]
	for every $r\geq 1$.
	Also choose $a_N\in G$ with $\mathfrak{s}_p(a_N)=K_N$, and set
	\[
		x=a_0+pa_1+\dots+p^Na_N.
	\]
	Fix $0\leq i\leq N$.
	By \autoref{lem:BasicPropertiesOfS},
	\[
		\mathfrak{s}_{p^{i+1}}(p^ia_i)=\mathfrak{s}_p(a_i)=K_i.
	\]
	If $j<i$, then
	\[
		\mathfrak{s}_{p^{i+1}}(p^ja_j)
		=
		\mathfrak{s}_{p^{i+1-j}}(a_j)
		<
		K_{j+1}
		\leq
		K_i,
	\]
	while, if $j>i$, then $\mathfrak{s}_{p^{i+1}}(p^ja_j)=\emptyset$.
	Thus, among the summands of $x$, the unique maximal $\mathfrak{s}_{p^{i+1}}$-value is $K_i$.
	Repeatedly applying \autoref{lem:BasicPropertiesOfS}(2) gives $\mathfrak{s}_{p^{i+1}}(x)=K_i$.
\end{proof}

\begin{thm}\label{thm:InfiniteSpine}
	If $\mathcal{S}_p$ is infinite for some prime $p$, then the dp-rank of $G$ is $\aleph_0$.
\end{thm}

\begin{proof}
	We work in a sufficiently saturated monster model of $\operatorname{Th}(G)$.
	Since $\mathcal{S}_p$ is an infinite interpretable linear order, compactness allows us to choose elements
	\[
		(b_{i,k})_{i,k<\omega}
	\]
	such that all the values
	\[
		H_{i,k}:=\mathfrak{s}_p(b_{i,k})
	\]
	are nonempty and
	\[
		H_{i,k}<H_{j,l}
		\quad\text{whenever}\quad
		i<j
		\ \text{or}\
		(i=j\ \text{and}\ k<l).
	\]
	Indeed, every finite part of these requirements merely asks for a finite chain of distinct nonempty elements of the infinite linear order $\mathcal{S}_p$.

	For a function $\eta:\omega\to\omega$, consider the partial type
	\[
		\Sigma_\eta(x)
		=
		\left\{
			\mathfrak{s}_{p^{i+1}}(x)=H_{i,\eta(i)}
			\;\middle|\;
			i<\omega
		\right\}.
	\]
	This type is finitely satisfiable.
	Indeed, given $N<\omega$, the values
	\[
		H_{0,\eta(0)}<H_{1,\eta(1)}<\dots<H_{N,\eta(N)}
	\]
	form a finite increasing chain, so \autoref{lem:FiniteCoding} realizes the first $N+1$ conditions of $\Sigma_\eta$.
	By saturation, there is $a_\eta\in G$ realizing $\Sigma_\eta$.

	For every $i<\omega$, let
	\[
		\varphi_i(x;y)
		\quad\text{be the formula}\quad
		\mathfrak{s}_{p^{i+1}}(x)=\mathfrak{s}_p(y).
	\]
	As noted above, this is an $L_{\mathrm{oag}}$-formula.
	For all $\eta:\omega\to\omega$, $i<\omega$, and $k<\omega$, we have
	\[
		\models\varphi_i(a_\eta;b_{i,k})
		\quad\Longleftrightarrow\quad
		H_{i,\eta(i)}=H_{i,k}
		\quad\Longleftrightarrow\quad
		\eta(i)=k.
	\]
	Hence, these formulas and parameters form an ICT-pattern of depth $\aleph_0$, and therefore the dp-rank of $G$ is at least $\aleph_0$.

	Every ordered abelian group has NIP \cite{Gurevich1984TheOrd}.
	Since $L_{\mathrm{oag}}$ is countable, this gives dp-rank at most $\aleph_0$.
	Consequently, the dp-rank of $G$ is exactly $\aleph_0$.
\end{proof}

\section{The lower bound in the finite-spine case}

We now assume that $G$ is nontrivial and has finite spines.
For a prime $p$ and $H\in\mathcal{S}_p\setminus\{\emptyset\}$, recall that
\[
	H^{[p]}=\bigcap_{K\Subset G,\,K>H}(K+pG)
\]
and
\[
	R_H^{\mathfrak{s}_p}
	=
	H^{[p]}/(H+pG)
\]
by \cite[Definition 2.10 and Lemma 2.14(1)]{Okura2026DisOrd}.
We also use $\mathcal{S}_{p^r}=\mathcal{S}_p$ for every $r\geq 1$, as established in \cite[Lemma 2.9(3)]{Okura2026DisOrd}.
We put
\[
	\mathcal{S}_p^\infty
	=
	\left\{
		H\in\mathcal{S}_p\setminus\{\emptyset\}
		\;\middle|\;
		\left|R_H^{\mathfrak{s}_p}\right|=\infty
	\right\}.
\]

\begin{lem}\label{lem:RepresentativesOfInfiniteRib}
	Let $p$ be a prime and $H\in\mathcal{S}_p^\infty$.
	There is a sequence $(a^k)_{k<\omega}$ such that
	\[
		\mathfrak{s}_p(a^k)
		=
		\mathfrak{s}_p(a^k-a^l)
		=
		H
	\]
	whenever $k\neq l$.
	If $H<K$ for some $K\in\mathcal{S}_p\setminus\{\emptyset\}$, the sequence may moreover be chosen inside $K$.
\end{lem}

\begin{proof}
	Choose elements $(u^k)_{k<\omega}$ whose images in
	\[
		R_H^{\mathfrak{s}_p}=H^{[p]}/(H+pG)
	\]
	are nonzero and pairwise distinct.
	Thus, neither $u^k$ nor $u^k-u^l$ belongs to $H+pG$, while all these elements belong to $H^{[p]}$.
	By \cite[Lemma 2.14(1)]{Okura2026DisOrd}, this says precisely that
	\[
		\mathfrak{s}_p(u^k)
		=
		\mathfrak{s}_p(u^k-u^l)
		=
		H
	\]
	for $k\neq l$.

	Suppose that $H<K$.
	Since $u^k\in H^{[p]}\subseteq K+pG$, we can write $u^k=a^k+pc_k$ with $a^k\in K$.
	Replacing $u^k$ by $a^k$ does not change its image in $R_H^{\mathfrak{s}_p}$.
	The images therefore remain nonzero and pairwise distinct, so the required equalities remain true.
\end{proof}

\begin{lem}[Simultaneous coding for one prime]\label{lem:SimultaneousCodingOnePrime}
	Let $p$ be a prime and let
	\[
		H_0<H_1<\dots<H_{n-1}
	\]
	be elements of $\mathcal{S}_p^\infty$.
	For each $i<n$, there is a sequence $(a_i^k)_{k<\omega}$ as in \autoref{lem:RepresentativesOfInfiniteRib}, and these sequences can be chosen so that $a_i^k\in H_{i+1}$ whenever $i<n-1$.
	For $\zeta:n\to\omega$, set
	\[
		b_\zeta
		=
		a_0^{\zeta(0)}
		+pa_1^{\zeta(1)}
		+\dots
		+p^{n-1}a_{n-1}^{\zeta(n-1)}.
	\]
	Then, for every $i<n$,
	\begin{enumerate}
		\item $\mathfrak{s}_{p^{i+1}}(b_\zeta)=H_i$;
		\item the image of $b_\zeta$ in $R_{H_i}^{\mathfrak{s}_{p^{i+1}}}$ is the same as that of $p^ia_i^{\zeta(i)}$;
		\item for every $k<\omega$,
		\[
			b_\zeta-p^ia_i^k\in H_i+p^{i+1}G
			\quad\Longleftrightarrow\quad
			k=\zeta(i).
		\]
	\end{enumerate}
\end{lem}

\begin{proof}
	The required sequences are supplied by \autoref{lem:RepresentativesOfInfiniteRib}.
	For $j<n-1$, the fact that $a_j^k\in H_{j+1}$ implies
	\[
		\mathfrak{s}_{p^r}(a_j^k)<H_{j+1}
	\]
	for every $r\geq 1$: otherwise, $\mathfrak{s}_{p^r}(a_j^k)\geq H_{j+1}$ would contain $a_j^k$, contrary to the definition of $\mathfrak{s}_{p^r}$.
	The calculation in the proof of \autoref{lem:FiniteCoding} now gives
	\[
		\mathfrak{s}_{p^{i+1}}(b_\zeta)=H_i,
	\]
	which proves (1).

	For $j<i$, we have $a_j^{\zeta(j)}\in H_{j+1}\subseteq H_i$, whereas for $j>i$, the term $p^ja_j^{\zeta(j)}$ belongs to $p^{i+1}G$.
	Hence,
	\[
		b_\zeta-p^ia_i^{\zeta(i)}
		\in
		H_i+p^{i+1}G.
	\]
	Also, by \autoref{lem:BasicPropertiesOfS},
	\[
		\mathfrak{s}_{p^{i+1}}(p^ia_i^{\zeta(i)})
		=
		\mathfrak{s}_p(a_i^{\zeta(i)})
		=
		H_i.
	\]
	Thus, both elements represent classes in
	\[
		R_{H_i}^{\mathfrak{s}_{p^{i+1}}}
		=
		H_i^{[p^{i+1}]}/(H_i+p^{i+1}G),
	\]
	and the congruence $b_\zeta - p^i a_i^{\zeta(i)} \in H_i + p^{i+1}G$ proves (2).

	If $k\neq\zeta(i)$, then
	\[
		p^i\left(a_i^{\zeta(i)}-a_i^k\right)
		\notin
		H_i+p^{i+1}G.
	\]
	Indeed, by the purity of the convex subgroup $H_i$, membership in the group on the right would imply
	\[
		a_i^{\zeta(i)}-a_i^k\in H_i+pG,
	\]
	contrary to \autoref{lem:RepresentativesOfInfiniteRib}.
	Together with the preceding congruence, this proves (3).
\end{proof}

The following standard consequences of compactness will be used without proof.
For the second assertion, finite satisfiability follows from the Chinese remainder theorem \cite[Lemma 2.1(2)]{Okura2026DisOrd}.

\begin{lem}\label{lem:CofinalDivisibleTranslations}
	Let $G$ be a sufficiently saturated monster model of the theory of a nontrivial ordered abelian group.
	\begin{enumerate}
		\item The type-definable subgroup
		\[
			D=\bigcap_{m\geq 2}mG
		\]
		is cofinal and coinitial in $G$.
		\item Let $P_0\subseteq\mathbb{P}$ and $(a_p)_{p\in P_0}$ be a sequence from $G$.
		There exists $a\in G$ such that
		\[
			a-a_p\in p^rG
		\]
		for every $p\in P_0$ and $r\geq 1$.
	\end{enumerate}
\end{lem}

\begin{thm}\label{thm:LowerBoundFiniteSpines}
	Let $G$ be a nontrivial ordered abelian group with finite spines.
	Then
	\[
		\operatorname{dp\text{-}rk}(G)
		\geq
		1+\sum_{p\in\mathbb{P}}\left|\mathcal{S}_p^\infty\right|.
	\]
	If the sum on the right is infinite, then the dp-rank of $G$ is $\aleph_0$.
\end{thm}

\begin{proof}
	We work in a sufficiently saturated monster model.
	For each prime $p$, put $n_p=\left|\mathcal{S}_p^\infty\right|$ and, when $n_p>0$, enumerate
	\[
		\mathcal{S}_p^\infty
		=
		\{H_{p,0}<\dots<H_{p,n_p-1}\}.
	\]
	The number $n_p$ is finite because $G$ has finite spines.
	Let
	\[
		\mathcal{I}
		=
		\{(p,i)\mid p\in\mathbb{P},\ 0\leq i<n_p\}.
	\]
	Thus,
	\[
		|\mathcal{I}|
		=
		\sum_{p\in\mathbb{P}}\left|\mathcal{S}_p^\infty\right|,
	\]
	and $\mathcal{I}$ is finite or countable.

	For each $p$ with $n_p>0$, choose sequences $(a_{p,i}^k)_{k<\omega}$ as in \autoref{lem:SimultaneousCodingOnePrime}, and denote the corresponding coding elements by $b_{p,\zeta}$ for $\zeta:n_p\to\omega$.
	For $(p,i)\in\mathcal{I}$, let
	\[
		\varphi_{p,i}(x;y)
		\quad\text{be the formula}\quad
		x-y\in\mathfrak{s}_{p^{i+1}}(y)+p^{i+1}G.
	\]
	This is an $L_{\mathrm{oag}}$-formula by the uniform definability of $\mathfrak{s}_{p^{i+1}}$.
	We use the parameter sequence
	\[
		\left(p^ia_{p,i}^k\right)_{k<\omega}
	\]
	in the row indexed by $(p,i)$.
	By \autoref{lem:BasicPropertiesOfS},
	\[
		\mathfrak{s}_{p^{i+1}}(p^ia_{p,i}^k)
		=
		H_{p,i}.
	\]

	Fix a function $\eta:\mathcal{I}\to\omega$ and consider
	\[
		q_\eta(x)
		=
		\left\{
			\varphi_{p,i}
			\left(x;p^ia_{p,i}^k\right)^{\eta(p,i)=k}
			\;\middle|\;
			(p,i)\in\mathcal{I},\ k<\omega
		\right\}.
	\]
	Here, the superscript means that the formula is negated when $\eta(p,i)\neq k$.
	For each $p$ with $n_p>0$, let
	\[
		\eta_p:n_p\to\omega,
		\qquad
		\eta_p(i)=\eta(p,i).
	\]
	By \autoref{lem:CofinalDivisibleTranslations}(2), there is $b_\eta\in G$ such that
	\[
		b_\eta-b_{p,\eta_p}\in p^rG
	\]
	for every $p$ with $n_p>0$ and $r\geq 1$.
	Part (3) of \autoref{lem:SimultaneousCodingOnePrime} then shows that $b_\eta\models q_\eta$.

	It remains to add the row arising from the order.
	Since the monster is sufficiently saturated, there exists $c>0$ such that
	\[
		-c<b_\eta<c
	\]
	for every $\eta:\mathcal{I}\to\omega$.
	By \autoref{lem:CofinalDivisibleTranslations}, choose
	\[
		t_0<t_1<\dots
	\]
	from $D$ so that $t_{k+1}-t_k>2c$ for every $k<\omega$.
	Set
	\[
		I_k=(t_k-c,t_k+c).
	\]
	These intervals are pairwise disjoint and satisfy $I_0<I_1<\dots$.
	We take
	\[
		\psi(x;y,z):=(y<x<z)
	\]
	for the additional row, with parameter sequence
	\[
		(t_k-c,t_k+c)_{k<\omega}.
	\]

	Let $\nu:\{\ast\}\mathbin{\dot\cup}\mathcal{I}\to\omega$, and put $\eta=\nu|_{\mathcal{I}}$.
	The element
	\[
		a_\nu=b_\eta+t_{\nu(\ast)}
	\]
	belongs to $I_{\nu(\ast)}$.
	Moreover, $t_{\nu(\ast)}\in p^{i+1}G$ for every $(p,i)\in\mathcal{I}$, so translation by $t_{\nu(\ast)}$ does not change the truth of any formula $\varphi_{p,i}$.
	It follows that
	\[
		\models\psi(a_\nu;t_k-c,t_k+c)
		\quad\Longleftrightarrow\quad
		\nu(\ast)=k
	\]
	and
	\[
		\models
		\varphi_{p,i}
		\left(a_\nu;p^ia_{p,i}^k\right)
		\quad\Longleftrightarrow\quad
		\nu(p,i)=k.
	\]
	We have therefore constructed an ICT-pattern of depth
	\[
		1+|\mathcal{I}|
		=
		1+\sum_{p\in\mathbb{P}}\left|\mathcal{S}_p^\infty\right|,
	\]
	as required.
	If the sum is infinite, the ICT-pattern has depth $\aleph_0$.
	The NIP upper bound recalled in the proof of \autoref{thm:InfiniteSpine} gives $\operatorname{dp\text{-}rk}(G)\leq\aleph_0$, so equality holds in this case.
\end{proof}

\addcontentsline{toc}{section}{References}
\begin{sloppypar}
	\printbibliography
\end{sloppypar}

\end{document}